\documentclass[12pt,a4paper,oneside]{report}
\usepackage[margin=25mm,headheight=20pt]{geometry}
\usepackage{fontspec}
\setmainfont{THSarabunPSK}[Path=assets/,Extension=.ttf,UprightFont=*-Regular,BoldFont=*-Bold,ItalicFont=*-Italic,BoldItalicFont=*-BoldItalic]
\setsansfont{THSarabunPSK}[Path=assets/,Extension=.ttf,UprightFont=*-Regular,BoldFont=*-Bold]
\setmonofont{lmmono10-regular.otf}[Scale=MatchLowercase]
% TH SarabunPSK: body 16pt, titles 18pt, small (header, diagram labels) 14pt
% sizes are Word/PDF points (bp): 16bp = 16.06 TeX pt
\makeatletter
\renewcommand\normalsize{\@setfontsize\normalsize{16.06}{20}%
  \abovedisplayskip 10\p@ \@plus3\p@ \@minus6\p@ \belowdisplayskip \abovedisplayskip
  \abovedisplayshortskip \z@ \@plus3\p@ \belowdisplayshortskip 6\p@ \@plus3\p@ \@minus3\p@}
\renewcommand\small{\@setfontsize\small{14.05}{17}}
\makeatother
\newcommand{\titlesize}{\fontsize{18.07}{22}\selectfont}
\normalsize
\XeTeXlinebreaklocale "th"
\XeTeXlinebreakskip=0pt plus 1pt
\usepackage{pifont,capt-of,graphicx,booktabs,tabularx,longtable,array,amsmath,enumitem,xcolor,tikz,fancyhdr,titlesec,float}
\usetikzlibrary{arrows.meta,positioning,shapes.geometric}
\usepackage[unicode,colorlinks=true,linkcolor=black,urlcolor=black,citecolor=black]{hyperref}
\usepackage{bookmark}
\definecolor{navy}{HTML}{000000}
\definecolor{pale}{HTML}{F2F2F2}
\renewcommand{\contentsname}{สารบัญ}
\renewcommand{\listtablename}{สารบัญตาราง}
\renewcommand{\listfigurename}{สารบัญภาพ}
\renewcommand{\tablename}{ตารางที่}
\renewcommand{\figurename}{ภาพที่}
\renewcommand{\bibname}{บรรณานุกรม}
\renewcommand{\chaptername}{บทที่}
\renewcommand{\appendixname}{ภาคผนวก}
\titleformat{\chapter}[display]{\bfseries\color{navy}\titlesize}{\chaptername\ \thechapter}{5pt}{\titlesize}
\titlespacing*{\chapter}{0pt}{0pt}{22pt}
\titleformat{\section}{\normalsize\bfseries\color{navy}}{\thesection}{0.6em}{}
\titleformat{\subsection}{\normalsize\bfseries}{\thesubsection}{0.6em}{}
\setlength{\parindent}{1.3em}
\setlength{\parskip}{5pt}
\linespread{1}
\setlength{\emergencystretch}{3em}
\tolerance=3000
\setlist{itemsep=3pt,topsep=5pt,leftmargin=1.6em}
\renewcommand{\arraystretch}{1.35}
\newcolumntype{Y}{>{\raggedright\arraybackslash}X}
\pagestyle{fancy}
\fancyhf{}
\fancyhead[L]{\small HENJAI | เห็นใจ}
\fancyhead[R]{\small ทีม Rushhour}
\fancyfoot[C]{\normalsize\thepage}
\renewcommand{\headrulewidth}{0.3pt}
\fancypagestyle{plain}{\fancyhf{}\fancyfoot[C]{\normalsize\thepage}\renewcommand{\headrulewidth}{0pt}}
\DeclareRobustCommand{\file}[1]{\nolinkurl{#1}}
\newcommand{\note}[1]{\par\noindent\colorbox{pale}{\parbox{\dimexpr\linewidth-2\fboxsep}{#1}}\par}
\newcommand{\frontchapter}[1]{\chapter*{#1}\addcontentsline{toc}{chapter}{#1}}
\newcommand{\projectcover}[1]{%
\hypersetup{pageanchor=false}\begin{titlepage}\centering\titlesize\bfseries
HENJAI\par
\vspace{4mm}เห็นใจ\par
\vspace{9mm}#1\par
\vspace{7mm}โครงการผู้ช่วยบรรยายภาพด้วยปัญญาประดิษฐ์\\ที่เข้าใจอารมณ์และบริบทสำหรับผู้พิการทางสายตา\par
\vspace{3mm}Emotion \& Context-Aware AI Vision Assistant\par
\vfill
ผู้จัดทำ\par\vspace{4mm}
นายกิตติพงศ์ สมสอาด\\
นายชินพัฒน์ ขุนเจริญ\\
นายเรข์ณพัศ สุวรรณปริสุทธิ์\par
\vspace{4mm}ทีม Rushhour\\นิสิตชั้นปีที่ 2 หลักสูตรวิศวกรรมคอมพิวเตอร์และเทคโนโลยีดิจิทัล (CEDT)\\ภาควิชาวิศวกรรมคอมพิวเตอร์\\คณะวิศวกรรมศาสตร์ จุฬาลงกรณ์มหาวิทยาลัย\par
\end{titlepage}\hypersetup{pageanchor=true}}
